\section{Exact precision computation and polynomial compilation}\label{sec:compiler}
Encode a signed integer by
\[
 \operatorname{enc}(n)=\begin{cases}2n,&n\ge0,\\-2n-1,&n<0.\end{cases}
\]
The five natural inputs decode in the fixed order $(n,\kappa,a,l,j)$.
The decision predicate rejects unless
$n\notin\Cset$, $\kappa\ne0$, $a,l\ge1$, and $j=na$.
For nonzero multiples of $P_0$, $z=-x/y$ is well-defined. Set
\[
 z([l]P_0)=u_l+v_ls,\qquad\epsilon_l=|u_l|+|v_l|,
 \qquad r(a,j)=\tfrac12\Tr_{K_0/\Q}\frac{z([j]P_0)}{z([a]P_0)}.
\]
With the fixed constant $A_{\rm ht}^{(0)}=H_6(S_{\rm ht})$, its final
tests are
\[
 |r(a,j)-n|<\epsilon_l,\qquad
 4\epsilon_l<4^{-A_{\rm ht}^{(0)}|\kappa F(n)|^{24}}.
\]
These are strict tests on exactly computable rationals and integers.

The reference program S21 uses exact quadratic-field arithmetic and capped
integer operations. The ideal register intermediate representation S22
computes the same predicate on every five-natural tuple. Its independent
all-input argument covers normalized signed arithmetic, origin and
zero-ordinate cases, equal-abscissa addition, and every early rejection.
Nine exact polynomial identities verify the nonexceptional projective
formulas and preservation of $Y^2=X^3+AXZ^4$. Scalar loops decrease by
halving; the precision loop divides by four; factorial iteration is bounded.
This is a proof about the ideal intermediate representation, rather than
an assertion that a budgeted simulator always finishes.

For example, if $\epsilon=U/V>0$ with natural $U,V$ and integer $B\ge0$,
\[
 4\epsilon<4^{-B}\quad\Longleftrightarrow\quad
 4^B\le\left\lfloor\frac{V-1}{4U}\right\rfloor.
\]
This permits a finite cap. For positive $L$, $\sigma(x)=\min(x,L+1)$
commutes with saturated addition and multiplication. Powers and factorials
can be computed with this saturation. Capping an exponent is safe for
base at least two because an exponent above $L$ already makes the value
exceed $L$. Base zero, base one and exponent zero are separate cases.
This proves equality with the fixed-height expression without requiring
its full evaluation when the available precision is too small.

\subsection{The natural polynomial and the power interface}
The saved precision polynomial has $1\,305\,811$ natural variables,
the first five being inputs. Its SHA256 is
\begin{center}\small\ttfamily
f3e163f5c60b972c98b48647274f956dd7d78428e27ddc0543574a5f18120b8c.
\end{center}
Source S23 specializes a fixed universal natural halting polynomial at
the code of a wrapper. The wrapper inverts four Cantor pairs and then
runs the unchanged precision machine. Its pairing equations are
\[
 2t_i=(c_i+t_{i+1})(c_i+t_{i+1}+1)+2t_{i+1}\quad(0\le i<4),
 \qquad t_4=c_4.
\]
Their triangular intervals prove unique natural unpairing. The new
polynomial sums the square of the old halting polynomial and the four
squared differences above.

The universal compiler packs a finite run of a fixed $m$-state
increment/decrement/branch program into base-$B$ integers. Acceptance
means reaching the unique final halt with all registers zero; the fixed
interpreter includes a finite cleanup stage. Choose
\[
 c>0,\quad K=2^c,\quad B=4K>m,\quad T=B^q,\quad
 F_q=1+B+\cdots+B^q,\quad D=(K-1)F_q,
\]
and require every input to be below $K$. For natural integers, say
$X$ is masked by $Y$ when every set binary bit of $X$ occurs in $Y$.
Require each register history $R_j$ to be masked by $D$ and satisfy
$R_j<T$. Each state history $S_i$ is masked by $F_q$, and
\[
 \sum_{i=0}^{m-1}S_i=F_q,\qquad S_{\rm halt}=T.
\]
Thus register digits are below $K$, terminal register digits are zero,
and state digits are zero or one. Since $B>m$, the state sum has no
carries and selects exactly one state at each time.

The nonzero mask $N_j$ is determined by
\[
 K N_j=(R_j+D)\mathbin{\mathrm{AND}}(K F_q).
\]
Indeed each digit $r+K-1<2K<B$ has its $K$-bit set exactly when
$r>0$. For a decrement/branch instruction at state $i$ testing
register $j$, impose
\[
 V_i=S_i\mathbin{\mathrm{AND}}N_j,\qquad S_i=V_i+W_i.
\]
These are the nonzero and zero branch histories. Form the packed
increment indicator $A_j$ from states incrementing $j$, the decrement
indicator $C_j$ from the appropriate $V_i$, and each incoming-state
indicator $I_i$ from the states or branches directed to $i$.
They have zero-one digits, since there is just one active instruction
and branch at each time. The transition equations are
\[
 R_j+BC_j=\operatorname{initial}_j+BR_j+BA_j,\qquad
 S_i=\mathbf1_{i=0}+BI_i.
\]
Every register coefficient on either side is at most $K<B$;
uniqueness of base-$B$ digits gives the initial conditions and each
exact transition. The possible extra final slot vanishes because
halt has no outgoing instruction and the terminal registers are zero.
For $q=0$ the same equations require the initial state already to be
halt and the inputs to vanish. Conversely, any finite accepting run
fits after choosing $c$ sufficiently large, and supplies all these
masks and equations.

\subsection{The literal alpha and power formulas}
Define $U_n=\alpha_b(n)$ by
\[
 U_0=0,\qquad U_1=1,\qquad U_{n+2}=bU_{n+1}-U_n\quad(n\ge0).
\]
For $b\ge3$ this recurrence gives
$0\le U_n<U_{n+1}$ and $U_n\ge n$. For $n\ge1$, the matrix
\[
 M=\begin{pmatrix}b&-1\\1&0\end{pmatrix},\qquad
 M^n=\begin{pmatrix}U_{n+1}&-U_n\\U_n&-U_{n-1}\end{pmatrix}
\]
has determinant one. Thus adjacent terms are coprime and
$U_n^2-bU_nU_{n-1}+U_{n-1}^2=1$.
Parameter congruence follows directly from the recurrence:
$b\equiv b'\pmod d$ implies
$\alpha_b(n)\equiv\alpha_{b'}(n)\pmod d$.
In particular $\alpha_2(n)=n$. This is parameter congruence,
not merely congruence between indices.

Every natural solution of $x^2+y^2=1+bxy$, $b>2$, is an unordered
adjacent pair. To prove this, order it as $X>Y$. If $Y=0$, $X=1$.
Otherwise $Z=bY-X$ satisfies $XZ=Y^2-1$, hence $0\le Z<Y$,
and $(Y,Z)$ satisfies the same equation. Descent reaches $(1,0)$;
reversing the steps generates exactly the adjacent pairs.
No ordering assumption on the original pair is needed.

For $k>0$, reduction of $M^k$ modulo $U_k$ is the invertible scalar
$U_{k+1}I$. Dividing an index by $k$ and using strict growth proves
$U_k\mid U_m$ if and only if $k\mid m$. More precisely, for $m>0$,
\[
                    U_k^2\mid U_m\quad\Longleftrightarrow\quad
                         kU_k\mid m.                           \tag{C.1}
\]
For $k=0$ both sides are false. For $k>0$ write $m=kq$ and expand
$M^{kq}=(U_kM-U_{k-1}I)^q$ modulo $U_k^2$. Its lower-left entry is
$qU_k(-U_{k-1})^{q-1}$, so coprimality proves (C.1), including $U_k=1$.

For $m>0$ put $v=U_{m+1}-U_{m-1}=\operatorname{tr}(M^m)$.
Cayley--Hamilton gives $M^{2m}\equiv-I\pmod v$.
Extend the recurrence by $U_{-j}=-U_j$. Every $n\ge0$ has
$n=2lm\pm j$ with $0\le j\le m$, and
$U_n\equiv\pm U_j\pmod v$. If $b\ge4$,
$v=bU_m-2U_{m-1}>2U_m$.

The actual alpha formula has natural witnesses $r,s,t,u,v,w,x,y$
and the following fifteen conditions:
\[
\begin{gathered}
 b>3,\quad u^2+t^2=1+but,\quad s^2+r^2=1+bsr,\quad r<s,\quad
 u^2\mid s,\quad v+2r=bs,\\
 w\equiv b\pmod v,\quad w\equiv2\pmod u,\quad w>2,\quad
 x^2+y^2=1+wxy,\\
 2X<u,\quad2X<v,\quad x\equiv X\pmod v,\quad
 2N<u,\quad x\equiv N\pmod u.
\end{gathered}                                                   \tag{C.2}
\]
A strict inequality $a<b$ is $b=a+1+d$; divisibility is $b=ad$;
congruence is $A+Mu=B+Mv$, with fresh natural witnesses.
At zero modulus the last equation says equality, so there is no
exception hidden in this replacement.

For soundness, the unordered Pell result gives $u=U_k$, $k>0$,
and $x=\alpha_w(n)$. The ordered pair gives $s=U_m,r=U_{m-1}$,
$m>0$. By (C.1), $u\mid m$. Parameter congruence gives
$x\equiv U_n\pmod v$ and $x\equiv n\pmod u$.
Reflect $n$ to $j\in[0,m]$. Since $2X<v$ and $2U_j<v$,
the congruence $X\equiv\pm U_j\pmod v$ forces $X=U_j$,
with the negative case forcing both zero.
Then $j\le X$, $2j<u$ and $2N<u$ force $N=j$ from
$N\equiv\pm j\pmod u$. Thus $X=U_N$.

Conversely choose an odd $u=U_k>2U_N$. Such terms are unbounded,
since adjacent terms are coprime and the sequence grows.
Set $t=U_{k+1}$, $m=ku$, $s=U_m$, $r=U_{m-1}$ and $v=bs-2r$.
Equation (C.1) gives $u^2\mid s$, and the required strict bounds hold.
The moduli $u,v$ are coprime: a common divisor is odd, divides $s$
and $2r$, hence divides both adjacent terms. CRT gives
$w\equiv2\pmod u$, $w\equiv b\pmod v$ with $w>2$.
Take $x=\alpha_w(N)$, $y=\alpha_w(N+1)$.
Parameter congruence verifies every remaining condition.
This includes $N=X=0$.

The literal power formula separates $Q=0,R=0,P=1$ and
$Q=0,R>0,P=0$. For $Q\ge1$ it imposes
\[
 a_0=\alpha_{Q+4}(R+1),\quad b=a_0+Q^2+2,\quad
 m+Q^2+1=bQ,\quad P<m,
\]
\[
 a_1=\alpha_b(R),\quad a_2=\alpha_b(R+1),\quad
 P+ba_1\equiv Qa_1+a_2\pmod m.                                  \tag{C.3}
\]
Induction gives $a_0\ge(Q+3)^R$, so $m>Q^R$ and $m>0$.
Let $E_R=Q U_R-bU_R+U_{R+1}$. The recurrence and
$bQ\equiv Q^2+1\pmod m$ give $E_0=1$ and
$E_{R+1}\equiv QE_R\pmod m$. Thus (C.3) says $P\equiv Q^R\pmod m$.
Both values lie in $[0,m)$, forcing equality. Conversely equality
supplies every witness. This includes $R=0$, $Q=1$ and $0^0=1$.

These formulas make the binomial and bit operations polynomial.
For $B=2^{n+1}$, the equation
\[
 (1+B)^n=B^k(Bq+C)+r,\quad C<B,\quad r<B^k
\]
selects exactly $\binom nk$, since all its base-$B$ coefficients
are below $B$, including zero for $k>n$.
The identity
$(1+t)^n=\prod_{i:\,\text{bit}_i(n)=1}(1+t^{2^i})$ over $\F_2$
makes odd binomial coefficients equivalent to bit inclusion.
Bit disjointness of $a,b$ is equivalent to inclusion of $a$ in $a+b$:
at the least shared bit there is no incoming carry and that bit
is cleared. Finally $x=z+a,y=z+b$ with $z,a,b$ pairwise bit-disjoint
says exactly $z=x\mathbin{\rm AND}y$.

\subsection{Interpreter, lowering and five-input specialization}
The universal program numbering is $e+1=2^w(2D+1)$.
At $w=0$ it halts. Otherwise the four base-$2^w$ digits at positions
$4p,\ldots,4p+3$ give the opcode, register and two targets at state $p$.
Opcode one increments, opcode two decrements or takes its zero branch;
every other opcode halts. Missing digits are zero, so decoding is total.

The interpreter stores a finite tape as
$0$ or $\operatorname{cons}(h,t)=2^h(2t+1)$.
Repeated halving decodes a nonempty head and tail; the tail is smaller
than the code. To update a register, save the finite reversed prefix,
replace the entry and restore the prefix. Absent entries are zero.
The reserved jump register remains zero. Copies restore their source;
aliased addition snapshots its operand before overwriting it.
Division maintains $x=qy+r$ and subtracts only when $r\ge y$;
the interpreter's divisors are at least two.
These invariants prove that one interpreter loop executes exactly one
decoded instruction, with finite work. A simulated halt enters cleanup
and clears all registers. An infinite simulation never enters cleanup.
Thus its all-zero halt is equivalent to virtual halting, without
requiring the virtual final tape itself to be zero.

The precision IR uses normalized signed pairs for integer arithmetic
and pairs over the basis $1,\sqrt2$ for field arithmetic.
Norm zero implies the field element is zero: after clearing
denominators, a primitive solution of $a^2=2b^2$ would have both
entries even. Every division is guarded by that exact test.
The Jacobian point formulas satisfy the curve equation and the affine
chord or tangent formulas after multiplying denominators.
Equal abscissas split into doubling and opposite-point branches;
zero ordinate doubles to the identity. The fixed projective point
is the same affine $P_0$. These are polynomial identities, independent
of the generator theorem.

The scalar loop halves its remaining count. The precision loop divides
its remainder by four. The capped arithmetic computes exactly
$\min(x,L+1)$; zero bases, unit bases and zero exponents are separate.
For bases at least two an exponent exceeding $L$ already exceeds the
cap. Thus the ideal IR terminates and agrees with the stated rational
predicate on all five natural inputs, including every early rejection.

Each primitive lowering uses fresh auxiliary registers, snapshots
overlapping sources, and preserves the live-register map at the end
of each finite block. Saturating decrement gives natural subtraction,
and repeated addition gives multiplication. All variable divisions
are guarded. Induction on completed loop iterations proves the same
claim for loops. Acceptance enters cleanup; rejection enters a
permanent self-loop. Hence the primitive machine halts exactly on
acceptance, although the ideal decision program is total.

For $\operatorname{pair}(a,b)=(a+b)(a+b+1)/2+b$, unpairing maintains
$\text{remainder}=\text{input}-s(s+1)/2$ and subtracts $s+1$
while $s<\text{remainder}$. At exit
$0\le\text{remainder}\le s$, so $b=\text{remainder}$, $a=s-b$.
Four repetitions put the five codes into registers zero through four.
Fresh scratch lies above the old register range and is cleared;
old work registers remain zero. Relocating every target by the prefix
length therefore preserves the precision machine.

The polynomial specialization fixes the universal program input to
this wrapper code, maps the old packed input to variable five, and
maps every old variable $i\ge2$ and private-block start by $i\mapsto i+7$.
The four pairing variables occupy five through eight, disjoint from
the five inputs. The appended four squared equations force their
unique Cantor values. Thus, for the exact saved polynomial $C$,
\[
 \forall\mathbf c\in\N^5,\qquad
 \exists\mathbf z\in\N^{1305806}\ C(\mathbf c,\mathbf z)=0
 \quad\Longleftrightarrow\quad
        \operatorname{accepts\_naturals}(\mathbf c)=\mathrm{True}. \tag{C.4}
\]
The natural compiler proof has no height, restricted-converse or
rational-model premise. Its use of the height expression concerns
only computation of that fixed integer expression.

\subsection{Translation to point codes}
Apply (C.4) using the exact signed index model. Induction on the actual
source DAG gives each node the code of its signed integer value.
Natural guards constrain source variables, not intermediate sums or
products. Raw addition receives canonical codes; globally typed
multiplication justifies the natural guards. Fixed generator equations
and scalar templates supply constants.

The power-value recipe has three signed code inputs, seventy-nine
private natural witnesses, and an output code for the value of the
power-template polynomial. That value is not independently required
to vanish. Its saved root equality is $[17751579,15]$.
The precision parent has $54056$ explicit variables and $15845$
disjoint private blocks of seventy-nine variables. No parent variable
node accesses a private coordinate. Replacing each block by its
private power-value call therefore preserves the full existential
assignment. Only the final parent root is equated to the identity.

At each primitive formula, equality is a difference, conjunction is
a sum of squares, and disjunction is a product. A topological induction,
including the globally conjoined fixed equations, proves equality of
polynomial values under substitution, not merely equisatisfiability
of an unrelated formula. This is the required link to the saved AST.

\section{A variable-shift relation for ordinary values}\label{sec:bridge}
We prove the central conditional deduction in full. Assume the exact
index model, the compiled precision relation, the height theorem and
$\mathsf{PC}$. Capital letters below denote five-coordinate point codes;
their lower-case indices are ordinary integers. The external variable
$t$ is any rational number.

\subsection{Approximation for every fixed shift}
Define $\Close(N,\mathcal K,t)$ by choosing $A,L,J$, linking all five
signed codes to precision inputs in the order $(n,\kappa,a,l,j)$, and
imposing $|t-r(a,j)|<\epsilon_l$ on those same coordinates. Put
\[
 B(n,\kappa)=H_6(S_{\rm ht})|\kappa F(n)|^{24}.
\]
The precision contract gives
\begin{align}
 \Close(\iota(n),\iota(\kappa),t)&\Longrightarrow
                  |t-n|<e^{-B(n,\kappa)},\label{eq:close-sound}\\
 \Close(\iota(n),\iota(\kappa),n)&\quad\text{for every }
                  n\notin\Cset,\ \kappa\ne0.\label{eq:close-complete}
\end{align}
The first follows from $|t-n|<2\epsilon_l<4^{-B}/2<e^{-B}$.
For the second, positive multiples of $P_0$ approach the identity
simultaneously in the compact group $E(\mathbf R)\times E(\mathbf R)$.
This follows, for instance, by partitioning the compact closure of the
cyclic subgroup into arbitrarily small neighborhoods and subtracting
two sufficiently close multiples. If the resulting positive differences
were bounded, a fixed nonzero multiple would be the identity; infinite
order excludes this. Choose $l$ first to make the positive $\epsilon_l$
small enough. Since the derivative of $[n]$ at the identity is $n$,
$r(a,na)\to n$ along positive multiples $a$ approaching it. Choose such
an $a$ and put $j=na$. This works for every already fixed nonzero shift
$\kappa$, regardless of its size.

\subsection{The six local tests}
For $N=\iota(n)$, $\mathcal K=\iota(\kappa)$, form the code for
$\kappa n$, choose $H=\iota(h)$, $U=\iota(s_0)$, and put
$V_0=H-\mathcal K$, $V_1=U-\iota(\kappa n)$.
Require $\kappa\ne0$, $H,V_0\ne O$. Define $w(Q)=1/x(Q)$ for
$Q\ne O$, $w(O)=0$. For $t\notin\Cset$ use
\begin{align*}
 D_1&=w(H),&D_2&=w(U),&D_3&=\sigma w(V_0),&D_4&=\sigma w(V_1),\\
 D_5&=z(U)/z(H)-t,&D_6&=\sigma(z(V_1)/z(V_0))-t.
\end{align*}
Here $\sigma$ is conjugation in $K_0$. Choose $\eta,q\in\Q$ and impose
\begin{align*}
 \eta^2&=1,&R_0&=\eta2730^6F(t),&a_0&=1-R_0q^2\ne0,\\
 R&=R_0a_0^2,&C_j&=D_j/F(t),&W&=\prod_{j=1}^6(1+RC_j^2)\ne0,\\
 \Theta&=\Tr_{K_0/\Q}(W^{-1})\ne0,&B_{\rm pole}&=R/\Theta^2,
 &\Phi(B_{\rm pole})&.
\end{align*}
Call this relation $\Test(N,\mathcal K,t)$. All quotients are guarded
polynomially. There is one common denominator $F(t)$, one product $W$,
one trace and one pole instance.

\begin{lemma}[Divisibility]\label{lem:test}
If $\Test(\iota(n),\iota(\kappa),t)$ holds, then
$M_{S_{\rm ht}}(t)\mid\kappa F(n)$.
\end{lemma}
\begin{proof}
Let $p\mid M_{S_{\rm ht}}(t)$. Exactly one $c\in\Cset$ has
$\delta=v_p(t-c)>0$ odd. The contact differences are units, so
$v_p(F(t))=\delta$; the field is unramified at $p$. Put $b=v_p(a_0)$.
If $q=0$, then $b=0$. Otherwise $\delta+2v_p(q)$ is odd and
nonzero, so $b=\min(0,\delta+2v_p(q))\le0$. In either case
$\alpha=v_p(R)=\delta+2b$ is odd. At each place above $p$,
\[
 v(1+RC_j^2)=\min(0,\alpha+2(v(D_j)-\delta)).
\]
Let $\lambda=\max_{v\mid p}v(W)\le0$. Then
$v_p(\Theta)\ge-\lambda\ge0$. The pole input has odd valuation
$\alpha-2v_p(\Theta)$, which is positive by Lemma~\ref{lem:pole}.
Therefore $\alpha>0$ and $\lambda> -\alpha/2$.

At a place attaining $\lambda$, each nonpositive factor valuation is
at least $\lambda$. Consequently
\[
 v(D_j)\ge\delta-3\alpha/4=\delta/4-3b/2>0
 \quad(1\le j\le6).
\]
Integral valuations are at least one, at the same place for all tests.
The first four put $H,U$ in one formal group and $V_0,V_1$ in the
conjugate formal group. At the identity we set both $z(O)=0$ and
$\log(O)=0$; thus the cases $U=O$ and $V_1=O$ obey the same ratio
identities below. Their denominators are nonzero by the guards on
$H$ and $V_0$. For a nonzero point in either formal group,
\[
 \log(Q)=z(Q)(1+e_Q),\qquad v(e_Q)\ge1.
\]
The invariant differential is integral, and integration loses at most
$v_p(r)$ in degree $r$; for $r\ge2$,
$(r-1)v(z)-v_p(r)\ge1$.
The relation $[s_0]H=[h]U$ gives
$\log(U)/\log(H)=s_0/h$. The fifth test makes this integral and
congruent to $t$ modulo $p$. Likewise the sixth gives
$(s_0-\kappa n)/(h-\kappa)\equiv t$.
Multiplying by the integer denominators and subtracting yields
$v_p(\kappa(n-t))\ge1$. Since $t\equiv c$, we get
$p\mid\kappa(n-c)$, hence $p\mid\kappa F(n)$. The squarefree
product divides as asserted. No unit assumption on $\kappa$ was used.
\end{proof}

\begin{lemma}[One suitable shift]\label{lem:test-complete}
Assume $\mathsf{PC}$. For every $n\in\Z\setminus\Cset$ there is
$\kappa\ne0$ with $\Test(\iota(n),\iota(\kappa),n)$.
\end{lemma}
\begin{proof}
Choose the sign so $R_0=\eta2730^6F(n)>0$ and let $S'$ be its finite
prime support. Each sufficiently small local formal neighborhood has
finite index in the compact local group. A positive common multiple
of the quotient orders gives $\kappa$ with $[\kappa]P_0$ in every
chosen neighborhood. Take
\[
 H=[2\kappa]P_0,\quad U=[2\kappa n]P_0,\quad
 V_0=[\kappa]P_0,\quad V_1=[\kappa n]P_0.
\]
The formal parameter and the derivative of $[n]$ make each $C_j$
integral above $S'$ after sufficiently shrinking the neighborhoods.
Infinite order preserves the nonzero guards.

Choose rational $q$ integral at $S'$ and arbitrarily close to
$1/\sqrt{R_0}$ in $\mathbf R$, without equality. Denominators
$R_0m+1$ with nearest integer numerators give such choices, adjusting
by one if equality occurs. At $S'$, $a_0$ is a unit and
$W\equiv1$, $\Theta\equiv2$ modulo $p^{v_p(R_0)}$.
Integral trace justifies this also at the ramified place two.
The pole input has valuation at least four at two and at least six
at the other primes of $S_0$. Outside $S'$, its valuation is
$2v_p(a_0)-2v_p(\Theta)$ and is even, so its reduced denominator is
a square. In both real embeddings, $R$ tends to zero positively,
$W\to1$ and $\Theta\to2$. All guards hold and
$0<B_{\rm pole}<1/2$ for a sufficiently close choice.
Hypothesis~\ref{hyp:PC} supplies the pole witnesses.
\end{proof}

\subsection{Exact ordinary value}
Define $\Val(N,t)$ by membership of $N$ and the disjunction
\[
 \bigvee_{c\in\Cset}(N=\iota(c)\wedge t=c)
 \quad\vee\quad
 \exists\mathcal K\,
 (\Close(N,\mathcal K,t)\wedge\Test(N,\mathcal K,t)).
\]
The same five coordinates $\mathcal K$ are used in both calls.
In an accepting main branch $n\notin\Cset$, $\kappa\ne0$, and
Lemma~\ref{lem:test} and the height estimate imply
\[
 h(t)\le H_6(S_{\rm ht})|\kappa F(n)|^{24}=B(n,\kappa).
\]
For $t=a/b\ne n$ in reduced form, $|a-nb|\ge1$, hence
\[
 |t-n|\ge1/b\ge e^{-h(t)}\ge e^{-B(n,\kappa)},
\]
contradicting~\eqref{eq:close-sound}. Thus $t=n$.
Conversely choose $\kappa$ by Lemma~\ref{lem:test-complete}, then use
\eqref{eq:close-complete} for that same shift. Private witness blocks
are disjoint, so the two calls can be satisfied simultaneously.
The six remaining integers have exactly their displayed branches.

Therefore $\Val(\iota(n),t)$ is exactly the graph $t=n$.
Existentially quantifying $N$ gives ordinary $\Z$ in $\Q$.
S01 and S16 multiply the exception and main-branch polynomials and
conjoin membership by a sum of squares. Under $\mathsf B$ this is
the saved root $P_*$. This proves the zero-set part of
Theorem~\ref{thm:conditional}.
